\documentclass[12px]{article}

\title{Lezione 8 Metodi Numerici}
\date{2026-03-24}
\author{Federico De Sisti}

\input{../../setup.tex}

\begin{document}
	\maketitle
	\newpage
	Prendiamo in considerazione
	\[
	\begin{cases}
		\dot y(t) = f(t,y(t)) \ \ \ t\in (0,T)\\
		y(0) = y_0
	\end{cases} \Rightarrow  y_{n+1} = y_n + h\Phi(t_n,y_n) + h\tau_{n+1}(h)
	.\] 
	\[
		\begin{cases}
			\bar u_{n+1} = \bar u_n + h \Phi(t_n,U\bar u_n) + \xi_{n+1} = \bar u_n + h[\Phi(t_n,\bar u_n) + \frac {\xi_{n+1}} {h} ]\\
		\bar u_0 = y_0 + \xi_0
		\end{cases}
	.\] 
	dove $\xi_{n+1}$ sono gli errori dovuti all'aritmetica finita.
	Scriviamo ora
	\[
	\begin{split}
		|y_{n+1} - \bar u_{n+1}| &= |y_{n+1} - u_{n+1} + u_{n+1} - \bar u_{n+1}|\\
					 &\leq |y_{n+1} - u_{n+1}| + |u_{n+1} - \bar u_{n+1}|
	\end{split}
	.\] 
	dove $u_{n+1}$ è la soluzione teorica (senza contare dell'errore floating point\\
	\[
		|y_{n+1} - u_{n+1}| \substack{FE\\\leq} \frac{e^{LT}-1}L = Ch^q
	.\] 
	che vale in tenerale per metodi consistenti di ordine $q$\\ 
	dalla zero stabilità:
	\[
		|w_n|\leq (|w_0| + h \sum |\delta_{k+1}|)e^{nh\Lambda}
	.\] 
	$\delta_0 = \xi_0 \ \ \delta_n = \frac{\xi_n}{h} \ \  z_n = \bar u_n$
	\[
	\begin{split}
		\Rightarrow |u_n - \bar u_n| &\leq \left(|\xi_0 - h \sum^{n-1}_{k=0}|\frac{\xi_{n+1}}h| \right) e^{nh\Lambda}\\
					     &\leq (|\xi| + h |\frac \xi h| n)e^{T\Lambda}\\
					     &\leq |\xi| \left( 1 + \frac T h \right)e^{T\Lambda}\\
	\end{split}
	.\] 
	dove abbiamo utilizzato $\xi = \sup_n |\xi_n|$\\
	Otteniamo quindi
	\[
		|y_{n+1} - \bar u_{n+1}|\leq Ch^q + |\xi| \left(1 + \frac Th \right)e^{T\Lambda}
	.\] 
 Lo schema va a convergenza ma l'errore dato da accumulo floatinc point esplode, c'è un errore di consistenza del metodo utilizzato,\\
 Cerchiamo un bilanciamento tra i due termini di errore $h^*$ tale che se\\
  \begin{itemize}
	  \item  $h > h^*$ l'errore di troncamento domina\\
	  \item  $h < h^*$ l'errore di arrotondamento domina
 \end{itemize}
 \textbf{Stabilità Asintotica per ODE}\\
 Descrive il comportamento delle soluzioni per tempi $T$ grandi $(T \rightarrow +\infty)$ \\
 Una ODE  è asintoticamente stabile se è stabile in senso Lyapunov $\dorall T > 0$ \\
 se è stabile in senso Lyapunov $\forall T> 0$
  \[
 |y(t) - z(t)| < C\varepsilon \ \ \begin{cases}
 	\dot y = f(t,y)\\
	y(0) = y_0
 \end{cases} \ \ \begin{cases}
 	\dot z = f(t,\dot z) + \delta(t)\\
	z(0) = y_0 + \delta_0
\end{cases} \ t\in [0,T]
 .\] 
 se $|\delta_0| < \varepsilon \ \ |\delta(t)|< \varepsilon \ \ \forall t$\\
 e se $|y(t) - z(t)| \xrightarrow{t \rightarrow +\infty} 0$ se $|\delta(t)| \xrightarrow {t \rightarrow +\infty}0$ \\
 \textbf{Stabilità Assoluta} (versione discreta)\\
 È una proprietà dello schema che descrive il comporamento della soluzoien numerica prodotta da un metodo a un passo per $t_n \rightarrow +\infty $ con $h$ fissato.\\
 \textbf{Osservazione}\\
 Nella zero-stabilità $t_n$ è fissato e $h \rightarrow 0 $, La definizione di stabilità assoluta si basa su una ode modello:
 \[
 \begin{cases}
 	\dot y (t) = \lambda y(t) \ \ \ t > 0 \\
	y(0) = 1 \ \ \ \ \ \lambda \in \C
\end{cases} \Rightarrow  y(t) = e^{\lambda t} = e^{Re(\lambda)t}e^{iIm(\lambda)t}
 .\]  
 \begin{center}
 \includegraphics[scale=0.16]{stabilita.jpg}
 \end{center}
 Se $Re(\lambda) < 0 $ allora  $y(t) \xrightarrow{t \rightarrow +\infty} 0$
 \begin{defi}
 	Un metodo numerico si dice assolutametne stabile se appplicato al problema lineare modello 
	\[
	\begin{cases}
		\dot y = \lambda y\\
		y(0) = 1
	\end{cases} \ \ \lambda\in \C \ \ Re(\lambda) < 0
	.\] 
	Produce una soluzione numerica $u_n$ tale che $|u_n| \xrightarrow{n \rightarrow +\infty}0$
 \end{defi}
 \textbf{Esempio} Eulero in avanti\\

 \[\begin{cases}
		\dot y = \lambda y\\
		y(0) = 1
	\end{cases}\ \ \begin{cases}
	u_{n+1} = u_n + h\lambda u_n = u_n ( 1 + h \lambda )\\
	u_0 = 1
	\end{cases}\] 
	$u_{n+1} = (1 + h\lambda) u_n = (1 + h\lambda)(u_{n-1}(1+h\lambda)) = (1 + h\lambda)^2u_{n-1} = (1+ h\lamda)^{n+1}u_0 = (1 + h\lambda)^{n+1}$\\
	$|u_n| = |(1+h\lambda)^n|\leq |1 + h\lambda|^n \xrightarrow{ n \rightarrow \infty} 0$\\
	se $|1 + h\lambda | < 1$ \\
	Ovvero se $h\lambda$ è nel cerchio unitario centrato in $(-1,0)$\\
	 \textbf{Caso oscillatore armonico}\\
	 $\ddot x = -x \ \ x\in \R$\\
	 $\dot x = v \ \ y = \icol {x\\v}$ \ \ $\dot y = \icol {0 & 1\\ -1 &0}\icol{x\\v} = \icol{v\\-x}$\\
	 $\dot v = -x \ \ \dot y = \icol {\dot x\\\dot v}$\\
	 $\dot y = \mu y$ \ \  $\mu = \icol {0 & 1\\-1 & 0} \rightarrow\lambda = \pm i$\\
	 A me no di un cambio di ocordinate la soluzione di questo problema corrisponde a risolvere
	 \[
	 \begin{cases}
	 	\dot z_1 = i z_1\\
		\dot z_2 = -i z_2
	 \end{cases}
	 .\] 
	 che è il problema modello con $\lambda = \pm i $ (non risolvibile con FE)\\
	  \textbf{Eulero implicito}\\
	  $
	  \begin{cases}
	  	\dot y = \lambda y\ \ \ \ Re(\lambda) < 0 \\ 
		y(0) = 1
	  \end{cases}
	  $\\
	  $u_{n+1} = u_n + hf(t_{n+1},u_{n+1}) = u_n + h\lambda u_{n+1}$\\
	  $u_{n=1} = u_n + h\lambda u_{n+1}$\\
	  $ \Rightarrow  (1-h\lambda)u_{n+1} = u_n$ \\
	  \[
		  u_{n+1} =\frac{ u_n}{(1-h\lambda)} = \frac{1}{(1-h\lambda)^{n+1}}\\
	  .\] 
	  \[
		  |u_{n+1}| = \left|\frac{1}{(1-h\lambda)^{n+1}}\right|\leq \left| \frac {1}{1-h\lambda}\right|^{n+1} \xrightarrow{n \rightarrow +\infty}0
	  .\] 
	  $ \Leftrightarrow \left| \frac 1{1-h\lambda} \right| <1$\\
	  $ \Leftrightarrow |1 - h\lambda| > 1$
\end{document}
